\chapter{Quarkonia und schwere Quarks}

\section{Ziel dieses Kapitels}

Dieses Kapitel behandelt gebundene Zustaende aus schweren Quarks und Antiquarks.
Solche Zustaende nennt man Quarkonia.

Quarkonia sind wichtig, weil sie eine Bruecke zwischen Spektroskopie, QCD,
Confinement und experimenteller Teilchenphysik bilden.

\begin{notebox}
Dieses Kapitel behandelt die Stichwoerter:
\begin{itemize}
    \item schwere Quarks,
    \item Charm-Quark,
    \item Bottom-Quark,
    \item Top-Quark,
    \item Quarkonium,
    \item Charmonium,
    \item Bottomonium,
    \item Positronium,
    \item Wasserstoffatom,
    \item Quark-Antiquark-Potential,
    \item gebundene Zustaende,
    \item Spektroskopie,
    \item Resonanzen,
    \item Zerfallsbreite,
    \item Lebensdauer,
    \item OZI-Unterdrueckung,
    \item offene und versteckte schwere Flavors.
\end{itemize}
\end{notebox}

\section{Schwere Quarks}

Im Standardmodell gibt es sechs Quark-Flavors:
\[
    u,\quad d,\quad s,\quad c,\quad b,\quad t.
\]

\begin{definitionbox}
Als \emph{schwere Quarks} bezeichnet man vor allem:
\[
    c,\qquad b,\qquad t.
\]
Das sind Charm-, Bottom- und Top-Quark.
\end{definitionbox}

\begin{conceptbox}
Die leichten Quarks
\[
    u,\quad d,\quad s
\]
bilden viele der bekannten Hadronenmultipletts.

Die schweren Quarks
\[
    c,\quad b,\quad t
\]
haben deutlich groessere Massen und fuehren zu eigenen Spektren schwerer Hadronen.
\end{conceptbox}

\begin{notebox}
Das Top-Quark ist so schwer und so kurzlebig, dass es nicht wie Charm- oder
Bottom-Quarks zu normalen gebundenen Quarkonium-Zustaenden hadronisiert.

Deshalb sind fuer Quarkonia vor allem Charmonium und Bottomonium wichtig.
\end{notebox}

\begin{examtipbox}
Fuer Quarkonia sind besonders wichtig:
\[
    c\overline{c}
    \quad\text{und}\quad
    b\overline{b}.
\]

Toponium
\[
    t\overline{t}
\]
bildet wegen der sehr kurzen Top-Lebensdauer keine normalen gebundenen Zustaende.
\end{examtipbox}

\section{Was ist ein Quarkonium?}

\begin{definitionbox}
Ein \emph{Quarkonium} ist ein gebundener Zustand aus einem schweren Quark und seinem
eigenen Antiquark:
\[
    Q\overline{Q}.
\]
Dabei steht $Q$ fuer ein schweres Quark.
\end{definitionbox}

\begin{conceptbox}
Die wichtigsten Quarkonia sind:
\[
    c\overline{c}
    \quad\Rightarrow\quad
    \text{Charmonium},
\]
\[
    b\overline{b}
    \quad\Rightarrow\quad
    \text{Bottomonium}.
\]
\end{conceptbox}

\begin{definitionbox}
\emph{Charmonium} ist ein gebundener Zustand aus Charm-Quark und Charm-Antiquark:
\[
    c\overline{c}.
\]

\emph{Bottomonium} ist ein gebundener Zustand aus Bottom-Quark und Bottom-Antiquark:
\[
    b\overline{b}.
\]
\end{definitionbox}

\begin{examtipbox}
Quarkonium bedeutet nicht irgendein Quark-Antiquark-Meson.

Der Begriff wird vor allem fuer schwere Systeme wie
\[
    c\overline{c}
    \quad\text{und}\quad
    b\overline{b}
\]
verwendet.
\end{examtipbox}

\section{Versteckter schwerer Flavor}

Quarkonia enthalten ein schweres Quark und das passende schwere Antiquark.
Dadurch ist die gesamte schwere Flavor-Quantenzahl null.

\begin{definitionbox}
Ein Hadron besitzt \emph{versteckten schweren Flavor}, wenn es ein schweres Quark
und das zugehoerige schwere Antiquark enthaelt.

Beispiele:
\[
    c\overline{c},
    \qquad
    b\overline{b}.
\]
\end{definitionbox}

\begin{conceptbox}
Beim Charmonium gilt:
\[
    C(c)=+1,
    \qquad
    C(\overline{c})=-1.
\]
Die gesamte Charm-Quantenzahl ist daher:
\[
    C_{\mathrm{ges}}=+1-1=0.
\]

Der Charm ist also im Zustand vorhanden, aber als Gesamtquantenzahl versteckt.
\end{conceptbox}

\begin{definitionbox}
Ein Hadron besitzt \emph{offenen schweren Flavor}, wenn es ein schweres Quark oder
Antiquark enthaelt, ohne dass derselbe Flavor durch den Partner kompensiert wird.
\end{definitionbox}

\begin{examplebox}
Beispiel fuer versteckten Charm:
\[
    J/\psi = c\overline{c}.
\]

Beispiel fuer offenen Charm:
\[
    D^0 = c\overline{u}.
\]

Beim $D^0$ ist die gesamte Charm-Quantenzahl nicht null.
\end{examplebox}

\section{Vergleich mit Wasserstoff und Positronium}

Quarkonia erinnern in manchen Aspekten an bekannte gebundene Zweiteilchensysteme.

\begin{center}
\begin{tabular}{|p{3.5cm}|p{3.4cm}|p{3.8cm}|p{3.5cm}|}
\hline
System & Bestandteile & Wechselwirkung & Potential \\
\hline
Wasserstoffatom & $p e^-$ & elektromagnetisch & Coulomb-Potential \\
\hline
Positronium & $e^- e^+$ & elektromagnetisch & Coulomb-Potential \\
\hline
Quarkonium & $Q\overline{Q}$ & stark & QCD-Potential \\
\hline
\end{tabular}
\end{center}

\begin{conceptbox}
Alle drei Systeme sind gebundene Zweiteilchensysteme.

Wasserstoff und Positronium werden durch elektromagnetische Wechselwirkung gebunden.

Quarkonia werden durch die starke Wechselwirkung gebunden.
\end{conceptbox}

\begin{examtipbox}
Die Analogie ist hilfreich, aber nicht perfekt.

Quarkonia sind keine rein coulombartigen Systeme.
Wegen Confinement besitzt das Quark-Antiquark-Potential bei grossen Abstaenden einen
linearen Anteil.
\end{examtipbox}

\section{Wasserstoffatom}

Das Wasserstoffatom besteht aus einem Proton und einem Elektron.

\begin{definitionbox}
Das \emph{Wasserstoffatom} ist ein gebundenes System aus einem Proton und einem
Elektron:
\[
    p e^-.
\]
\end{definitionbox}

\begin{conceptbox}
Die Bindung entsteht durch das elektromagnetische Coulomb-Potential:
\[
    V(r)
    =
    -\frac{1}{4\pi\varepsilon_0}
    \frac{e^2}{r}.
\]

Das Spektrum des Wasserstoffatoms besteht aus diskreten Energieniveaus.
\end{conceptbox}

\begin{notebox}
Das Proton ist viel schwerer als das Elektron.
Deshalb bewegt sich das Elektron naeherungsweise im Potential eines fast ruhenden
Protons.

Genauer verwendet man die reduzierte Masse des Zweiteilchensystems.
\end{notebox}

\section{Positronium}

Positronium besteht aus Elektron und Positron.

\begin{definitionbox}
\emph{Positronium} ist ein gebundener Zustand aus Elektron und Positron:
\[
    e^-e^+.
\]
\end{definitionbox}

\begin{conceptbox}
Positronium wird ebenfalls elektromagnetisch gebunden.
Das Potential ist coulombartig.

Der Unterschied zum Wasserstoffatom ist:
Elektron und Positron haben gleiche Masse.
Dadurch ist die reduzierte Masse anders als im Wasserstoffatom.
\end{conceptbox}

\begin{notebox}
Positronium ist instabil.
Elektron und Positron koennen annihilieren und in Photonen zerfallen.
\end{notebox}

\begin{examtipbox}
Wasserstoff:
\[
    p e^-.
\]

Positronium:
\[
    e^-e^+.
\]

Quarkonium:
\[
    Q\overline{Q}.
\]
\end{examtipbox}

\section{Quarkonium als stark gebundenes Zweiteilchensystem}

Quarkonium besteht aus einem schweren Quark und seinem Antiquark.

\begin{conceptbox}
Weil die Quarks schwer sind, kann man Quarkonia in erster Naeherung manchmal
nichtrelativistisch behandeln.

Das macht sie spektral aehnlich zu Atomphysik-Systemen:
Es gibt Grundzustaende, angeregte Zustaende und Uebergaenge zwischen Niveaus.
\end{conceptbox}

\begin{notebox}
Die Analogie zur Atomphysik ist besonders gut bei schweren Quarks.
Je schwerer die Quarks, desto langsamer bewegen sie sich im gebundenen Zustand
naeherungsweise relativ zueinander.
\end{notebox}

\begin{examtipbox}
Quarkonia sind wichtig, weil sie das QCD-Potential experimentell zugaenglich machen.
Ihre Spektren enthalten Information ueber die starke Wechselwirkung.
\end{examtipbox}

\section{Quark-Antiquark-Potential}

Das Quark-Antiquark-Potential beschreibt die Bindung zwischen Quark und Antiquark.

\begin{formulabox}
Eine haeufig verwendete qualitative Form ist:
\[
    V(r)
    =
    -\frac{a}{r}
    +
    \kappa r.
\]
Dabei beschreibt:
\begin{itemize}
    \item $-\frac{a}{r}$ den coulombartigen Anteil bei kleinen Abstaenden,
    \item $\kappa r$ den linearen Confinement-Anteil bei grossen Abstaenden.
\end{itemize}
\end{formulabox}

\begin{conceptbox}
Bei kleinen Abstaenden ist die starke Kopplung wegen asymptotischer Freiheit kleiner.
Das Potential besitzt dort einen coulombartigen Anteil.

Bei grossen Abstaenden wird der lineare Term wichtig.
Er beschreibt das Confinement.
\end{conceptbox}

\begin{examtipbox}
Quark-Antiquark-Potential:
\[
    \text{kleines } r:
    \quad
    -\frac{a}{r}.
\]

\[
    \text{grosses } r:
    \quad
    \kappa r.
\]

Der lineare Anteil verhindert freie Quarks.
\end{examtipbox}

\section{Spektroskopie von Quarkonia}

Quarkonia besitzen diskrete Energiezustaende.
Diese Zustaende werden durch Quantenzahlen beschrieben.

\begin{conceptbox}
Wichtige Quantenzahlen sind:
\begin{itemize}
    \item Hauptquantenzahl beziehungsweise radialer Anregungszustand,
    \item Bahndrehimpuls $L$,
    \item Spin $S$ des Quark-Antiquark-Systems,
    \item Gesamtdrehimpuls $J$,
    \item Paritaet $P$,
    \item C-Paritaet $C$.
\end{itemize}
\end{conceptbox}

\begin{notebox}
In der Spektroskopie verwendet man oft die Schreibweise:
\[
    n\,{}^{2S+1}L_J.
\]

Dabei ist:
\begin{itemize}
    \item $n$ die radiale Quantenzahl,
    \item $S$ der Gesamtspin von Quark und Antiquark,
    \item $L$ der Bahndrehimpuls,
    \item $J$ der Gesamtdrehimpuls.
\end{itemize}
\end{notebox}

\begin{examtipbox}
Die Schreibweise
\[
    {}^{2S+1}L_J
\]
kommt aus der Spektroskopie.

Sie ist nicht nur in der Atomphysik, sondern auch bei Quarkonia nuetzlich.
\end{examtipbox}

\section{Bahndrehimpuls-Notation}

\begin{notebox}
Fuer den Bahndrehimpuls verwendet man Buchstaben:
\[
    L=0 \Rightarrow S,
\]
\[
    L=1 \Rightarrow P,
\]
\[
    L=2 \Rightarrow D,
\]
\[
    L=3 \Rightarrow F.
\]
\end{notebox}

\begin{conceptbox}
Achtung:
Der Buchstabe $S$ in der Spektroskopie kann zwei verschiedene Dinge meinen.

In
\[
    {}^{2S+1}L_J
\]
steht $S$ im Exponenten fuer den Gesamtspin.

Der Buchstabe $S$ als Bahndrehimpulsbezeichnung bedeutet dagegen
\[
    L=0.
\]
\end{conceptbox}

\begin{examtipbox}
Beispiel:
\[
    {}^3S_1
\]
bedeutet:
\begin{itemize}
    \item Spin-Triplett: $S=1$,
    \item Bahndrehimpuls: $L=0$,
    \item Gesamtdrehimpuls: $J=1$.
\end{itemize}
\end{examtipbox}

\section{Paritaet und C-Paritaet von Quarkonia}

Quarkonia sind Mesonen aus Quark und Antiquark.
Daher kann man Paritaet und C-Paritaet aus $L$ und $S$ bestimmen.

\begin{formulabox}
Fuer ein Quark-Antiquark-System gilt:
\[
    P=(-1)^{L+1}.
\]
\end{formulabox}

\begin{formulabox}
Fuer ein neutrales Quarkonium gilt:
\[
    C=(-1)^{L+S}.
\]
\end{formulabox}

\begin{conceptbox}
Damit kann man fuer Quarkonium-Zustaende die Quantenzahlen
\[
    J^{PC}
\]
bestimmen.
Diese Schreibweise ist in der Mesonenspektroskopie sehr wichtig.
\end{conceptbox}

\begin{examplebox}
Betrachte einen Zustand
\[
    {}^3S_1.
\]
Dann gilt:
\[
    L=0,
    \qquad
    S=1,
    \qquad
    J=1.
\]

Paritaet:
\[
    P=(-1)^{0+1}=-1.
\]

C-Paritaet:
\[
    C=(-1)^{0+1}=-1.
\]

Also:
\[
    J^{PC}=1^{--}.
\]
\end{examplebox}

\begin{examplebox}
Betrachte einen Zustand
\[
    {}^1S_0.
\]
Dann gilt:
\[
    L=0,
    \qquad
    S=0,
    \qquad
    J=0.
\]

Paritaet:
\[
    P=(-1)^{0+1}=-1.
\]

C-Paritaet:
\[
    C=(-1)^{0+0}=+1.
\]

Also:
\[
    J^{PC}=0^{-+}.
\]
\end{examplebox}

\section{Charmonium}

Charmonium ist ein gebundener Zustand aus Charm-Quark und Charm-Antiquark.

\begin{definitionbox}
\emph{Charmonium} ist:
\[
    c\overline{c}.
\]
\end{definitionbox}

\begin{conceptbox}
Wichtige Charmonium-Zustaende sind:
\begin{itemize}
    \item $\eta_c$,
    \item $J/\psi$,
    \item $\psi(2S)$,
    \item $\chi_c$-Zustaende.
\end{itemize}
\end{conceptbox}

\begin{notebox}
Das $J/\psi$ ist besonders wichtig historisch und experimentell.
Es ist ein Vektor-Charmonium-Zustand mit Quantenzahlen
\[
    J^{PC}=1^{--}.
\]
\end{notebox}

\begin{examtipbox}
Charmonium:
\[
    c\overline{c}.
\]

Wichtigster Name:
\[
    J/\psi.
\]
\end{examtipbox}

\section{Bottomonium}

Bottomonium ist ein gebundener Zustand aus Bottom-Quark und Bottom-Antiquark.

\begin{definitionbox}
\emph{Bottomonium} ist:
\[
    b\overline{b}.
\]
\end{definitionbox}

\begin{conceptbox}
Wichtige Bottomonium-Zustaende sind:
\begin{itemize}
    \item $\eta_b$,
    \item $\Upsilon(1S)$,
    \item $\Upsilon(2S)$,
    \item $\Upsilon(3S)$,
    \item $\chi_b$-Zustaende.
\end{itemize}
\end{conceptbox}

\begin{notebox}
Bottomonium ist oft noch nichtrelativistischer als Charmonium, weil das Bottom-Quark
schwerer ist als das Charm-Quark.
Daher ist die Spektroskopie von Bottomonium besonders gut fuer Potentialmodelle
geeignet.
\end{notebox}

\begin{examtipbox}
Bottomonium:
\[
    b\overline{b}.
\]

Wichtigster Name:
\[
    \Upsilon.
\]
\end{examtipbox}

\section{Warum Quarkonia schmale Resonanzen sein koennen}

Manche Quarkonium-Zustaende sind relativ schmal.
Das bedeutet, sie haben eine kleine Zerfallsbreite und eine vergleichsweise lange
Lebensdauer.

\begin{definitionbox}
Die \emph{Zerfallsbreite} $\Gamma$ beschreibt die Breite einer Resonanz und haengt
mit ihrer Lebensdauer zusammen.
\end{definitionbox}

\begin{formulabox}
Qualitativ gilt:
\[
    \tau \sim \frac{\hbar}{\Gamma}.
\]
\end{formulabox}

\begin{conceptbox}
Eine kleine Breite bedeutet:
\[
    \Gamma \text{ klein}
    \quad\Rightarrow\quad
    \tau \text{ gross}.
\]

Eine grosse Breite bedeutet:
\[
    \Gamma \text{ gross}
    \quad\Rightarrow\quad
    \tau \text{ klein}.
\]
\end{conceptbox}

\begin{examtipbox}
Schmale Resonanz:
kleine Zerfallsbreite und relativ lange Lebensdauer.

Breite Resonanz:
grosse Zerfallsbreite und kurze Lebensdauer.
\end{examtipbox}

\section{Schwelle fuer offene schwere Flavors}

Ein Quarkonium-Zustand kann unterhalb oder oberhalb der Schwelle fuer offene schwere
Hadronen liegen.

\begin{definitionbox}
Die \emph{Schwelle fuer offenen Charm} ist die Energie, ab der ein Charmonium-Zustand
in Hadronen mit offenem Charm zerfallen kann, zum Beispiel:
\[
    D\overline{D}.
\]
\end{definitionbox}

\begin{conceptbox}
Wenn ein Charmonium-Zustand unterhalb der
\[
    D\overline{D}
\]
-Schwelle liegt, kann er nicht einfach stark in zwei offene Charm-Mesonen zerfallen.

Dann sind manche Zerfaelle unterdrueckt, und der Zustand kann relativ schmal sein.
\end{conceptbox}

\begin{definitionbox}
Analog ist die \emph{Schwelle fuer offenen Bottom} die Energie, ab der Bottomonium
in offene Bottom-Hadronen zerfallen kann, zum Beispiel:
\[
    B\overline{B}.
\]
\end{definitionbox}

\begin{examtipbox}
Unterhalb der offenen Flavor-Schwelle sind starke Zerfaelle in offene schwere
Hadronen kinematisch verboten.
Das kann Quarkonium-Zustaende schmal machen.
\end{examtipbox}

\section{OZI-Unterdrueckung}

Bestimmte Zerfaelle von Quarkonia sind unterdrueckt, obwohl sie nicht durch einfache
Quantenzahlen verboten sind.

\begin{definitionbox}
Die \emph{OZI-Regel} beschreibt die Unterdrueckung von Prozessen, bei denen
Quarklinien getrennt werden muessen und der Anfangsquarkinhalt nicht direkt in den
Endzustand durchlaeuft.
\end{definitionbox}

\begin{conceptbox}
Bei einem Quarkonium-Zustand
\[
    Q\overline{Q}
\]
muss das schwere Quark-Antiquark-Paar fuer manche Zerfaelle annihilieren.

Solche Zerfaelle koennen gegenueber offenen Flavor-Zerfaellen unterdrueckt sein.
\end{conceptbox}

\begin{notebox}
Die OZI-Regel ist eine qualitative Regel.
Sie hilft zu verstehen, warum manche Hadronenresonanzen schmaler sind als naiv
erwartet.
\end{notebox}

\section{Leptonische Zerfaelle von Vektor-Quarkonia}

Vektor-Quarkonia mit
\[
    J^{PC}=1^{--}
\]
koennen in ein Lepton-Antilepton-Paar zerfallen.

\begin{formulabox}
Beispiele:
\[
    J/\psi \rightarrow e^+e^-,
\]
\[
    \Upsilon \rightarrow \mu^+\mu^-.
\]
\end{formulabox}

\begin{conceptbox}
Solche Zerfaelle laufen ueber elektromagnetische Annihilation:
\[
    Q\overline{Q}
    \rightarrow
    \gamma^*
    \rightarrow
    \ell^+\ell^-.
\]

Dabei steht
\[
    \ell
\]
fuer ein geladenes Lepton:
\[
    e,\quad \mu,\quad \tau.
\]
\end{conceptbox}

\begin{examtipbox}
Vektor-Quarkonia mit
\[
    J^{PC}=1^{--}
\]
sind in
\[
    e^+e^-
\]
-Kollisionen besonders gut erzeugbar, weil sie die Quantenzahlen eines virtuellen
Photons besitzen.
\end{examtipbox}

\section{Erzeugung von Quarkonia}

Quarkonia koennen in verschiedenen Experimenten erzeugt werden.

\begin{conceptbox}
Moegliche Produktionsmechanismen:
\begin{itemize}
    \item Elektron-Positron-Annihilation,
    \item Proton-Proton-Kollisionen,
    \item Photoproduktion,
    \item Zerfaelle schwerer Hadronen.
\end{itemize}
\end{conceptbox}

\begin{examplebox}
In Elektron-Positron-Kollisionen:
\[
    e^+e^- \rightarrow \gamma^* \rightarrow J/\psi.
\]

Das funktioniert besonders gut fuer Zustaende mit
\[
    J^{PC}=1^{--}.
\]
\end{examplebox}

\begin{examplebox}
In Proton-Proton-Kollisionen koennen schwere Quark-Antiquark-Paare erzeugt werden:
\[
    pp \rightarrow c\overline{c}+X
\]
oder
\[
    pp \rightarrow b\overline{b}+X.
\]

Ein Teil dieser Paare kann zu Quarkonia hadronisieren.
\end{examplebox}

\begin{notebox}
LHCb ist ein LHC-Experiment, das besonders stark mit Beauty-Physik verbunden ist.
Der Name LHCb steht fuer LHC beauty.
\end{notebox}

\section{Zerfallskanaele und experimentelle Signaturen}

Quarkonia werden experimentell ueber ihre Zerfallsprodukte nachgewiesen.

\begin{conceptbox}
Typische Signaturen:
\begin{itemize}
    \item zwei Myonen:
    \[
        \mu^+\mu^-,
    \]
    \item zwei Elektronen:
    \[
        e^+e^-,
    \]
    \item Photonen aus Uebergaengen zwischen Quarkonium-Zustaenden,
    \item Hadronen aus starken Zerfaellen.
\end{itemize}
\end{conceptbox}

\begin{examtipbox}
Ein sehr typischer experimenteller Nachweis ist:
\[
    J/\psi \rightarrow \mu^+\mu^-.
\]

Man rekonstruiert die invariante Masse des Myonpaares.
Ein Peak bei der $J/\psi$-Masse zeigt das Resonanzsignal.
\end{examtipbox}

\section{Invariante Masse bei Quarkonia}

Quarkonia erscheinen experimentell als Peaks in invarianten-Masse-Verteilungen.

\begin{formulabox}
Fuer zwei Zerfallsprodukte mit Energien
\[
    E_1,\qquad E_2
\]
und Impulsen
\[
    \vec p_1,\qquad \vec p_2
\]
gilt:
\[
    M^2c^4
    =
    (E_1+E_2)^2
    -
    \left|\vec p_1+\vec p_2\right|^2c^2.
\]
\end{formulabox}

\begin{conceptbox}
Wenn die zwei Teilchen aus einem Quarkonium-Zerfall stammen, ist
\[
    M
\]
die rekonstruierte Masse des Quarkonium-Zustands.

Viele Ereignisse ergeben eine Verteilung.
Ein Quarkonium-Zustand erscheint als Peak.
\end{conceptbox}

\begin{examtipbox}
Resonanzsuche:
\begin{enumerate}
    \item Zerfallsprodukte messen.
    \item Invariante Masse berechnen.
    \item Nach Peaks suchen.
    \item Peakposition gibt Masse.
    \item Peakbreite haengt mit Lebensdauer und Detektoraufloesung zusammen.
\end{enumerate}
\end{examtipbox}

\section{Uebergaenge zwischen Quarkonium-Zustaenden}

Quarkonia koennen zwischen angeregten und niedrigeren Zustaenden uebergehen.

\begin{conceptbox}
Moegliche Uebergaenge:
\begin{itemize}
    \item Photonemission:
    \[
        Q\overline{Q}^{*}
        \rightarrow
        Q\overline{Q}
        +
        \gamma,
    \]
    \item hadronische Uebergaenge, zum Beispiel mit Pionen,
    \item Zerfall in Leptonpaare oder Hadronen.
\end{itemize}
\end{conceptbox}

\begin{notebox}
Solche Uebergaenge sind analog zu Spektrallinien in der Atomphysik, aber die
zugrunde liegende Wechselwirkung ist die starke Wechselwirkung.
Photonen koennen trotzdem bei elektromagnetischen Uebergaengen emittiert werden.
\end{notebox}

\section{Feinstruktur und Hyperfeinstruktur}

Quarkonium-Spektren zeigen Aufspaltungen zwischen Zustaenden mit verschiedenen
Spin- und Drehimpulsquantenzahlen.

\begin{definitionbox}
\emph{Feinstruktur} bezeichnet Aufspaltungen, die mit Bahndrehimpuls, Spin-Bahn-Kopplung
und Tensorwechselwirkungen zusammenhaengen.

\emph{Hyperfeinstruktur} bezeichnet Aufspaltungen zwischen verschiedenen
Spin-Kopplungen der Quarks.
\end{definitionbox}

\begin{examplebox}
S-Zustaende mit
\[
    L=0
\]
koennen Spin-Singulett oder Spin-Triplett sein.

Spin-Singulett:
\[
    {}^1S_0.
\]

Spin-Triplett:
\[
    {}^3S_1.
\]

Diese Zustaende haben unterschiedliche Massen.
\end{examplebox}

\begin{conceptbox}
Solche Aufspaltungen liefern Information ueber die spinabhaengigen Teile der starken
Wechselwirkung.
\end{conceptbox}

\section{Charmonium versus Bottomonium}

\begin{center}
\begin{tabular}{|p{4.0cm}|p{5.0cm}|p{5.0cm}|}
\hline
Eigenschaft & Charmonium & Bottomonium \\
\hline
Quarkinhalt & $c\overline{c}$ & $b\overline{b}$ \\
\hline
typischer Name & $J/\psi$, $\psi$ & $\Upsilon$ \\
\hline
Quarkmasse & schwer & noch schwerer \\
\hline
nichtrelativistische Naeherung & brauchbar & meist besser \\
\hline
offene Flavor-Schwelle & $D\overline{D}$ & $B\overline{B}$ \\
\hline
\end{tabular}
\end{center}

\begin{conceptbox}
Bottomonium ist wegen der groesseren Bottom-Quark-Masse oft naeher am
nichtrelativistischen Potentialmodell.

Charmonium ist leichter und relativistische beziehungsweise gekoppelte Kanal-Effekte
koennen staerker sein.
\end{conceptbox}

\section{Top-Quark und fehlendes Toponium}

\begin{conceptbox}
Das Top-Quark zerfaellt extrem schnell ueber die schwache Wechselwirkung:
\[
    t \rightarrow b + W^+.
\]

Es lebt so kurz, dass es normalerweise keine Zeit hat, einen gebundenen
Toponium-Zustand zu bilden.
\end{conceptbox}

\begin{notebox}
Das ist ein wichtiger Unterschied zwischen:
\[
    c,\quad b
\]
und
\[
    t.
\]

Charm- und Bottom-Quarks bilden Hadronen.
Top-Quarks zerfallen vorher.
\end{notebox}

\begin{examtipbox}
Charmonium:
ja.

Bottomonium:
ja.

Toponium:
keine normalen gebundenen Spektren, weil das Top-Quark zu schnell zerfaellt.
\end{examtipbox}

\section{Quarkonia und Confinement}

Quarkonia sind direkte Beispiele fuer Confinement.

\begin{conceptbox}
Ein Quarkonium-Zustand besteht aus einem farbigen Quark und einem farbigen Antiquark.
Das gesamte System ist aber farbneutral.

Wenn man Quark und Antiquark weit voneinander trennen will, waechst die Feldenergie.
Statt freier Quarks entstehen neue Hadronen.
\end{conceptbox}

\begin{formulabox}
Anschaulich:
\[
    Q\overline{Q}
    \rightarrow
    Q\overline{q}
    +
    q\overline{Q}
\]
wenn die Energie fuer ein leichtes
\[
    q\overline{q}
\]
-Paar reicht.
\end{formulabox}

\begin{conceptbox}
Oberhalb der offenen Flavor-Schwelle kann ein Quarkonium-Zustand daher in zwei
offene schwere Mesonen zerfallen.
\end{conceptbox}

\section{Exotische quarkoniumartige Zustaende}

Neben normalen Quarkonia gibt es auch quarkoniumartige Zustaende, die nicht einfach
als
\[
    Q\overline{Q}
\]
beschrieben werden.

\begin{conceptbox}
Moegliche Interpretationen exotischer quarkoniumartiger Zustaende:
\begin{itemize}
    \item Tetraquarks,
    \item Molekuele aus zwei Mesonen,
    \item hybride Zustaende mit gluonischer Anregung,
    \item gekoppelte Kanal-Effekte.
\end{itemize}
\end{conceptbox}

\begin{notebox}
Fuer die Grundvorlesung ist meist wichtiger:
Normales Quarkonium ist
\[
    Q\overline{Q}.
\]
Exotische Zustaende zeigen, dass die QCD reichhaltigere farbneutrale Strukturen
erlaubt.
\end{notebox}

\section{Vergleich: Atomphysik und Quarkoniumphysik}

\begin{center}
\begin{tabular}{|p{4.0cm}|p{5.0cm}|p{5.0cm}|}
\hline
Aspekt & Atomphysik & Quarkoniumphysik \\
\hline
Bindung & elektromagnetisch & stark \\
\hline
Austauschteilchen & Photon & Gluonen \\
\hline
Potential & Coulomb-artig & Coulomb-artig plus linear \\
\hline
freie Bestandteile & Elektron und Proton frei beobachtbar & Quarks nicht frei beobachtbar \\
\hline
Spektrum & diskrete Energieniveaus & diskrete Resonanzen \\
\hline
Uebergaenge & Photonenemission & Photonen und Hadronen moeglich \\
\hline
\end{tabular}
\end{center}

\begin{examtipbox}
Der entscheidende Unterschied:
Im Wasserstoffatom kann man Elektron und Proton getrennt beobachten.

Im Quarkonium kann man Quark und Antiquark wegen Confinement nicht frei isolieren.
\end{examtipbox}

\section{Mini-Zusammenfassung}

\begin{notebox}
Die wichtigsten Punkte dieses Kapitels sind:
\begin{itemize}
    \item Quarkonia sind gebundene Zustaende aus schwerem Quark und Antiquark:
    \[
        Q\overline{Q}.
    \]
    \item Charmonium ist:
    \[
        c\overline{c}.
    \]
    \item Bottomonium ist:
    \[
        b\overline{b}.
    \]
    \item Toponium bildet keine normalen gebundenen Zustaende, weil das Top-Quark zu
    schnell zerfaellt.
    \item Quarkonia sind analog zu Wasserstoff und Positronium, aber durch die starke
    Wechselwirkung gebunden.
    \item Das Quark-Antiquark-Potential hat qualitativ die Form:
    \[
        V(r)=-\frac{a}{r}+\kappa r.
    \]
    \item Bei kleinen Abstaenden ist der coulombartige Anteil wichtig.
    \item Bei grossen Abstaenden ist der lineare Confinement-Anteil wichtig.
    \item Quarkonium-Zustaende werden mit Quantenzahlen wie
    \[
        J^{PC}
    \]
    klassifiziert.
    \item Fuer Quark-Antiquark-Systeme gilt:
    \[
        P=(-1)^{L+1}.
    \]
    \item Fuer neutrale Quarkonia gilt:
    \[
        C=(-1)^{L+S}.
    \]
    \item Quarkonia erscheinen experimentell als Peaks in invarianten-Masse-Verteilungen.
    \item Die Zerfallsbreite haengt mit der Lebensdauer zusammen:
    \[
        \tau \sim \frac{\hbar}{\Gamma}.
    \]
\end{itemize}
\end{notebox}

\section{Pruefungscheck}

\begin{examtipbox}
Nach diesem Kapitel solltest du folgende Fragen beantworten koennen:
\begin{enumerate}
    \item Was sind schwere Quarks?
    \item Was ist ein Quarkonium?
    \item Was ist Charmonium?
    \item Was ist Bottomonium?
    \item Warum gibt es kein normales Toponium-Spektrum?
    \item Was bedeutet versteckter schwerer Flavor?
    \item Was bedeutet offener schwerer Flavor?
    \item Wie vergleicht man H-Atom, Positronium und Quarkonium?
    \item Was ist der wichtigste Unterschied zwischen Positronium und Quarkonium?
    \item Welche qualitative Form hat das Quark-Antiquark-Potential?
    \item Was bedeutet der lineare Term im Potential?
    \item Was bedeutet die Notation $n\,{}^{2S+1}L_J$?
    \item Was bedeuten $S$, $P$, $D$ in der Spektroskopie?
    \item Wie berechnet man Paritaet bei Quarkonia?
    \item Wie berechnet man C-Paritaet bei neutralen Quarkonia?
    \item Welche Quantenzahlen hat ein ${}^3S_1$-Zustand?
    \item Warum ist $J/\psi$ wichtig?
    \item Warum ist $\Upsilon$ wichtig?
    \item Was ist eine offene Charm- beziehungsweise Bottom-Schwelle?
    \item Warum koennen manche Quarkonia schmale Resonanzen sein?
    \item Wie weist man Quarkonia experimentell nach?
    \item Warum erscheinen Quarkonia als Peaks in invarianten-Masse-Verteilungen?
\end{enumerate}
\end{examtipbox}